% =========================
% chuong/02-co-so-ly-thuyet.tex
% MỤC 2. CƠ SỞ LÝ THUYẾT
% Tài liệu: ưu tiên từ 2021, KHÔNG dùng MDPI, Frontiers, Hindawi. DOI đã kiểm qua Crossref ngày 05/11/2026.
% =========================

\section{CƠ SỞ LÝ THUYẾT}

% ------------------------------------------------------------
\subsection{Tổng quan các khái niệm liên quan}

\subsubsection{Ứng dụng web trong hoạt động kinh doanh thương mại điện tử}

Trong thương mại điện tử, ứng dụng web có hai nhóm vai trò. Nhóm thứ nhất hướng tới khách hàng:
website bán hàng, trang sản phẩm, giỏ hàng, thanh toán. Nhóm thứ hai phục vụ vận hành nội bộ: các
ứng dụng web quản lý đơn hàng, chăm sóc khách hàng, báo cáo cho quản lý. Nhóm thứ hai không trực tiếp
tạo ra đơn hàng nhưng quyết định chất lượng dịch vụ sau bán, từ đó ảnh hưởng tới đánh giá, uy tín và
khả năng khách quay lại mua.

So với phần mềm cài trên máy, nền tảng web phù hợp cho nghiệp vụ vận hành ở bốn điểm: người dùng chỉ
cần trình duyệt, nhiều người làm việc cùng lúc trên cùng một dữ liệu, quyền truy cập được phân theo
vai trò, và doanh nghiệp cập nhật hệ thống một lần cho mọi người dùng. Hệ thống trong đồ án thuộc nhóm
thứ hai: một ứng dụng web nội bộ giúp bộ phận chăm sóc khách hàng biết nên liên hệ ai trước, ghi lại
đã làm gì, và giúp quản lý theo dõi chất lượng chăm sóc qua các con số.

\subsubsection{Sự hài lòng của khách hàng và đánh giá trực tuyến}

Sự hài lòng của khách hàng là mức độ trải nghiệm thực tế đáp ứng hoặc vượt kỳ vọng của khách
hàng đối với sản phẩm và dịch vụ. Trong thương mại điện tử, sự hài lòng được thể hiện rõ nhất qua
điểm đánh giá sau khi nhận hàng. Điểm đánh giá vừa là phản hồi cho doanh nghiệp, vừa là thông tin
công khai mà người mua sau tham khảo, nên một đánh giá thấp ảnh hưởng trực tiếp đến doanh số của
người bán \parencite{deshpande2023}. Tuy nhiên, chỉ một phần khách hàng chủ động viết đánh giá, và
những người viết thường có trải nghiệm rất tốt hoặc rất tệ \parencite{li2008}. Vì vậy, đánh giá là
tín hiệu có giá trị nhưng không đại diện đầy đủ cho toàn bộ khách hàng.

\subsubsection{Sự không hài lòng ẩn}

Trong đồ án này, \textit{sự không hài lòng ẩn} (hidden dissatisfaction) được hiểu là trường hợp
khách hàng chấm 1 hoặc 2 sao cho một đơn hàng mà các chỉ số vận hành không báo bất thường, cụ thể là
đơn đã được giao trước hoặc đúng ngày dự kiến. Những đơn này không xuất hiện trong các báo cáo theo
dõi giao trễ, nên doanh nghiệp chỉ biết khách không hài lòng khi đánh giá đã được đăng. Như đã trình
bày ở mục 1.1, nhóm này chiếm 66,3\% số đơn bị chấm thấp trong dữ liệu Olist.

\subsubsection{Khôi phục dịch vụ và chăm sóc khách hàng chủ động}

Khôi phục dịch vụ (service recovery) là các hành động doanh nghiệp thực hiện để khắc phục một trải
nghiệm không tốt của khách hàng. \textcite{vanvaerenbergh2019} mô tả khôi phục dịch vụ như một hành
trình gồm ba giai đoạn: trước khi khách phàn nàn, trong lúc xử lý và sau khi xử lý; doanh nghiệp có
thể can thiệp ngay từ giai đoạn đầu thay vì chờ khách lên tiếng. Nghiên cứu của
\textcite{nazifi2021} trong ngành hàng không cho thấy chủ động liên hệ khách trước khi sự cố được
khách cảm nhận giúp giảm phản hồi tiêu cực trên mạng rõ rệt so với xử lý bị động. Hệ thống trong đồ
án áp dụng tinh thần này: nhân viên liên hệ khách có nguy cơ không hài lòng trước khi họ để lại đánh
giá.

\subsubsection{Hệ thống hỗ trợ ra quyết định và vai trò của con người}

Hệ thống hỗ trợ ra quyết định (Decision Support System) là hệ thống thông tin giúp người dùng ra
quyết định dựa trên dữ liệu và mô hình, nhưng quyết định cuối cùng vẫn thuộc về con người. Đây là
điểm quan trọng khi hệ thống dùng kết quả của mô hình học máy. Tổng quan của \textcite{burton2020}
cho thấy người dùng thường e ngại thuật toán (algorithm aversion), nhất là
sau khi thấy thuật toán sai. \textcite{dietvorst2018} chỉ ra người dùng chấp nhận thuật toán chưa
hoàn hảo hơn nếu họ được điều chỉnh kết quả, dù chỉ một chút. Thí nghiệm của \textcite{bansal2021}
cho thấy việc kèm lời giải thích không tự động làm đội người và AI ra quyết định tốt hơn, vì người dùng
có thể tin theo cả những gợi ý sai. Từ các kết quả này, nhóm thiết kế hệ thống để điểm rủi ro chỉ là
căn cứ ưu tiên: nhân viên đọc yếu tố liên quan, tự liên hệ khách, tự ghi kết quả, và giao diện luôn
nhắc rằng điểm số không phải kết luận.

\subsubsection{Bảng điều khiển (dashboard)}

\textcite{sarikaya2019} phân loại dashboard theo mục đích sử dụng và chỉ ra dashboard vận hành cần
giúp người xem nhận ra việc phải làm, chứ không chỉ hiển thị số liệu. \textcite{bach2023} tổng hợp
các mẫu thiết kế dashboard, trong đó có nguyên tắc đặt chỉ số quan trọng nhất ở vị trí dễ thấy, nhóm
các chỉ số liên quan và cho phép đi từ con số tổng hợp tới dữ liệu chi tiết. Trang Tổng quan của hệ
thống áp dụng các nguyên tắc này: số đơn mức cao chờ giao được đặt đầu tiên, mỗi con số đều bấm được
để mở danh sách đơn đã lọc sẵn.

% ------------------------------------------------------------
\subsection{Các nghiên cứu và dự án liên quan}

Nhóm chia các nghiên cứu liên quan thành bốn hướng.

\textbf{Giao hàng và đánh giá của khách hàng.} \textcite{deshpande2023} dùng dữ liệu hơn 15 triệu
đơn trên Tmall và chứng minh hiệu quả giao hàng ảnh hưởng tới điểm đánh giá, từ đó ảnh hưởng tới
doanh số. \textcite{ravula2023} cho thấy giao trễ làm điểm giảm mạnh, còn giao sớm chỉ được cộng rất
ít. \textcite{brylla2022} chỉ ra giao nhanh bằng cách tách đơn có thể phản tác dụng. Các nghiên cứu
này giải thích vì sao giao đúng hạn là cần nhưng chưa đủ, nhưng không đưa ra công cụ để doanh nghiệp
xử lý từng đơn.

\textbf{Dự đoán sự hài lòng trên dữ liệu Olist.} \textcite{orman2026} so sánh các mô hình học sâu để
phân loại đơn hài lòng và không hài lòng. Hướng này tập trung vào độ chính xác của mô hình và dừng ở
nhãn dự đoán; kết quả chưa được đưa vào một quy trình mà nhân viên có thể dùng hằng ngày.

\textbf{Ưu tiên xử lý khi nguồn lực có hạn.} \textcite{ascarza2018} chứng minh khách hàng có rủi ro
cao nhất chưa chắc là nhóm phản hồi tốt nhất với việc chăm sóc. \textcite{lemmens2020} đề xuất chọn
khách để can thiệp theo lợi nhuận kỳ vọng thay vì chỉ theo xác suất rời bỏ. \textcite{li2026} cho thấy
xếp hạng bằng học máy không phải lúc nào cũng tốt hơn quy tắc đơn giản. Bài học cho hệ thống là danh
sách ưu tiên phải hiển thị thêm thông tin nghiệp vụ (SLA, thời gian chờ, người bán) để con người tự
điều chỉnh, và phải ghi lại kết quả xử lý để sau này đánh giá được hiệu quả.

\textbf{Can thiệp chủ động và cộng tác người--máy.} \textcite{nazifi2021} và
\textcite{vanvaerenbergh2019} ủng hộ can thiệp sớm; \textcite{dietvorst2018} và
\textcite{bansal2021} cho thấy người dùng cần quyền điều chỉnh và cần hiểu giới hạn của gợi ý.

Bảng~\ref{tab:nghien-cuu} tóm tắt các nghiên cứu chính và phần hệ thống kế thừa.

{\small
\renewcommand{\arraystretch}{1.25}
\begin{longtable}{|>{\raggedright\arraybackslash}p{2.4cm}|>{\raggedright\arraybackslash}p{4.1cm}|>{\raggedright\arraybackslash}p{3.4cm}|>{\raggedright\arraybackslash}p{4.3cm}|}
\caption{Các nghiên cứu liên quan và phần hệ thống kế thừa}\label{tab:nghien-cuu}\\
\hline
\textbf{Nghiên cứu} & \textbf{Kết quả chính} & \textbf{Hạn chế với bài toán} & \textbf{Hệ thống kế thừa} \\ \hline
\endfirsthead
\hline
\textbf{Nghiên cứu} & \textbf{Kết quả chính} & \textbf{Hạn chế với bài toán} & \textbf{Hệ thống kế thừa} \\ \hline
\endhead
\textcite{deshpande2023} & Giao hàng ảnh hưởng tới điểm đánh giá và doanh số & Không xét đơn giao đúng hạn mà vẫn bị chê & Lý do cần xử lý trước khi khách đánh giá \\ \hline
\textcite{ravula2023} & Giao trễ làm điểm giảm mạnh, giao sớm cộng ít & Chỉ xét thời gian giao & Không coi ``đúng hạn'' là ``đã ổn'' \\ \hline
\textcite{orman2026} & So sánh mô hình phân loại sự hài lòng trên Olist & Dừng ở nhãn, chưa có quy trình sử dụng & Kết quả mô hình là đầu vào, hệ thống lo phần sử dụng \\ \hline
\textcite{ascarza2018}; \textcite{lemmens2020} & Rủi ro cao chưa chắc nên ưu tiên & Bối cảnh rời bỏ dịch vụ & Ghi kết quả xử lý để đo hiệu quả về sau \\ \hline
\textcite{li2026} & Học máy không luôn hơn quy tắc đơn giản & Bài toán giao vận & Hiển thị SLA, thời gian chờ cạnh điểm số \\ \hline
\textcite{nazifi2021} & Chủ động liên hệ giảm phản hồi tiêu cực & Ngành hàng không & Quy trình liên hệ khách sớm có SLA \\ \hline
\textcite{dietvorst2018}; \textcite{bansal2021} & Người dùng cần quyền điều chỉnh, giải thích không tự đủ & Thí nghiệm có kiểm soát & Con người quyết định, có ghi chú phản hồi về mô hình \\ \hline
\textcite{bach2023}; \textcite{sarikaya2019} & Mẫu thiết kế dashboard vận hành & Không gắn với nghiệp vụ cụ thể & Bố cục trang Tổng quan, bấm từ số tới đơn \\ \hline
\end{longtable}
}

% ------------------------------------------------------------
\subsection{Lý thuyết và mô hình được áp dụng}

\subsubsection{Kiến trúc client--server và ứng dụng web nhiều trang}

Hệ thống theo kiến trúc client--server ba tầng. Tầng trình bày là các trang web chạy trên trình
duyệt; tầng xử lý là máy chủ API kiểm tra quyền và quy tắc nghiệp vụ; tầng dữ liệu là cơ sở dữ liệu
quan hệ. Trình duyệt và máy chủ trao đổi qua giao thức HTTP \parencite{rfc9110}. Giao diện được xây
theo mô hình nhiều trang (multi-page application): mỗi màn hình là một tệp HTML riêng, dùng chung
một tệp CSS và các tệp JavaScript; sau khi tải trang, JavaScript gọi API theo kiểu AJAX để lấy dữ liệu JSON rồi dựng
nội dung bằng DOM. Cách làm này không cần framework, đúng với các kiến thức HTML, CSS, JavaScript,
DOM và JSON của học phần.

\subsubsection{HTML, CSS và JavaScript}

\begin{itemize}
    \item \textbf{HTML5} mô tả cấu trúc trang: thanh điều hướng, bảng, biểu mẫu (form) với các ô
    \texttt{input}, \texttt{select}, \texttt{textarea}; thuộc tính \texttt{role} và \texttt{aria-*} giúp
    trình đọc màn hình hiểu các hộp thoại.
    \item \textbf{CSS3} định dạng giao diện. Hệ thống dùng biến CSS cho bảng màu và khoảng cách,
    Flexbox và Grid để dàn trang, và media query để đổi bố cục theo độ rộng màn hình (thiết kế
    đáp ứng - responsive design).
    \item \textbf{JavaScript} xử lý sự kiện người dùng, kiểm tra dữ liệu nhập trước khi gửi, gọi
    API theo kỹ thuật AJAX bằng đối tượng \texttt{XMLHttpRequest} \parencite{mdnxhr}, đọc và dựng
    dữ liệu JSON lên trang qua DOM (kể cả các hộp thoại xác nhận được tạo bằng DOM), và lưu một
    số thiết lập (ngôn ngữ, mã đăng nhập) ở \texttt{localStorage}.
\end{itemize}

\subsubsection{RESTful API và định dạng JSON}

REST là phong cách kiến trúc trong đó mỗi tài nguyên có một địa chỉ (URL) và được thao tác bằng các
phương thức HTTP chuẩn \parencite{fielding2000}. Hệ thống dùng \texttt{GET} để đọc, \texttt{POST}
để tạo bản ghi mới hoặc thực hiện một thao tác nghiệp vụ (nhận đơn, hoàn tất, mở lại),
\texttt{PATCH} để cập nhật thiết lập người dùng. Mã trạng thái HTTP cho biết kết quả: 200/201 khi
thành công, 401 khi chưa đăng nhập, 403 khi không đủ quyền, 404 khi không tìm thấy, 409 khi trạng
thái không cho phép (ví dụ hai người cùng nhận một đơn) và 422 khi dữ liệu gửi lên không hợp lệ.
Dữ liệu trao đổi ở dạng JSON, tên trường viết kiểu \texttt{snake\_case}, thời gian theo chuẩn ISO
8601 có múi giờ.

\subsubsection{Xác thực bằng JWT và lưu mật khẩu an toàn}

JSON Web Token (JWT) là chuỗi mã có chữ ký số, chứa thông tin người dùng và thời hạn hiệu lực
\parencite{rfc7519}. Sau khi đăng nhập đúng, máy chủ cấp JWT; trình duyệt gửi kèm mã này trong tiêu
đề \texttt{Authorization: Bearer} ở mỗi yêu cầu, máy chủ kiểm tra chữ ký để biết ai đang gọi và vai
trò gì. Mật khẩu không được lưu nguyên văn mà lưu dạng băm bằng thuật toán Argon2, là thuật toán được
OWASP khuyến nghị cho lưu trữ mật khẩu \parencite{owaspargon}.

\subsubsection{FastAPI và PostgreSQL}

FastAPI là framework Python để xây dựng API \parencite{fastapidocs}. Mỗi endpoint khai báo kiểu dữ
liệu vào và ra bằng thư viện Pydantic, nhờ đó dữ liệu sai kiểu bị từ chối tự động và tài liệu API
(OpenAPI) được sinh ra từ mã nguồn. PostgreSQL là hệ quản trị cơ sở dữ liệu quan hệ mã nguồn mở
\parencite{postgres16}. Ngoài khóa chính, khóa ngoại và ràng buộc \texttt{CHECK}, hệ thống dùng
trigger của PostgreSQL để buộc một số bảng chỉ được thêm mới (append-only), và dùng cột phiên bản
dòng (\texttt{row\_version}) để tránh hai người cùng sửa một case cùng lúc.

\subsubsection{Docker Compose và Nginx}

Docker đóng gói mỗi thành phần cùng môi trường chạy của nó thành một container; Docker Compose khai
báo và khởi động nhiều container cùng lúc bằng một tệp cấu hình \parencite{dockercompose}. Hệ thống
gồm bốn dịch vụ: cơ sở dữ liệu, bước tạo bảng (migration), máy chủ API và máy chủ web Nginx. Nginx
vừa phục vụ các trang tĩnh vừa chuyển các yêu cầu \texttt{/api/v1} sang máy chủ API, nên trình duyệt
chỉ làm việc với một địa chỉ duy nhất.

\subsubsection{Các mô hình phân tích và thiết kế hệ thống}

\begin{itemize}
    \item \textbf{BPMN 2.0} (Business Process Model and Notation) mô tả quy trình nghiệp vụ bằng sự
    kiện bắt đầu, kết thúc, công việc, cổng rẽ nhánh và các làn (lane) thể hiện ai làm việc gì
    \parencite{omg2014bpmn}.
    \item \textbf{Sơ đồ phân rã chức năng (FDD)} chia hệ thống thành các nhóm chức năng và chức năng
    con theo dạng cây.
    \item \textbf{Sơ đồ luồng dữ liệu (DFD)} mô tả dữ liệu đi từ tác nhân bên ngoài qua các tiến
    trình xử lý tới kho dữ liệu. Nhóm dùng ký hiệu Gane--Sarson với ba mức: sơ đồ ngữ cảnh, mức 0 và
    mức 1.
    \item \textbf{Sơ đồ Use Case} của UML mô tả tác nhân và các chức năng họ sử dụng, cùng các quan hệ
    «include», «extend» và kế thừa giữa tác nhân \parencite{omg2017uml}.
    \item \textbf{Sơ đồ thực thể - liên kết (ERD)} mô tả các bảng dữ liệu và quan hệ giữa chúng
    \parencite{chen1976}; nhóm vẽ theo ký hiệu chân quạ (crow's foot) để thể hiện bản số một - nhiều.
\end{itemize}

\subsubsection{Thiết kế giao diện đáp ứng và dễ tiếp cận}

Giao diện được thiết kế để dùng được trên nhiều kích thước màn hình bằng media query và bố cục linh
hoạt. Nhóm tham khảo hướng dẫn WCAG 2.2 \parencite{w3c2023wcag} ở ba điểm: độ tương phản giữa chữ và
nền, không dùng màu là cách duy nhất để truyền đạt thông tin (mức rủi ro luôn có chữ và biểu tượng đi
kèm màu), và mọi nút có nhãn đọc được.
